\section{Extension from a stratum and its residue link}\label{sec:links}

Proposition~\ref{prop:strata-adjunction} gives normal dlt pairs on the
strata and a single actual adjoint line along every incidence.  We now
show that a section on the floor of a stratum extends after it satisfies
at most one birational residue comparison.  The comparison comes from
two markings of a projective-line fiber on a Mori model.  In the next
section we will control the actions generated by these comparisons.

\subsection{Semiample lines on the normal strata}

We use a small model both to place lower-dimensional abundance inside its
explicit sheaf conventions and to run the later perturbed program.  Let
\(Z\) be a positive-dimensional stratum of dimension at most three.
Lower its remaining floor coefficients as in the preceding section,
obtaining an effective klt boundary \(B_Z^-\) with actual adjoint.
There is a projective small ordinary \(\Q\)-factorialization
\(q_Z:Z_0\to Z\) of this pair.  Here is a construction that also checks
the actual line used on it.

On a projective log resolution \(r:W\to Z\), keep the effective strict
boundary and assign to each exceptional prime a rational weight just below
one but strictly above its crepant coefficient.  This produces an
effective SNC klt boundary \(C_W\) with
\[
 K_W+C_W=r^*(K_Z+B_Z^-)+F,
\]
where \(F\) is effective, exceptional, and positive on every divisorial
exceptional component.  The equality is the meromorphic actual-line
identity.  The adjoint is relatively pseudo-effective.  Apply the
relative MMP for a projective morphism of compact analytic spaces in
dimension at most three \cite[Proposition~2.26]{DasHacon2024}.  Its source
is the ordinary \(\Q\)-factorial smooth \(W\), its pair is effective SNC
klt, and its morphism is a projective surjection of normal compact
analytic spaces.  On the
relative ordinary \(\Q\)-factorial minimal model over \(Z\), the transform
of \(F\) is exceptional and relatively nef, since the pulled-back adjoint
is relatively trivial.  Negativity makes that transform zero.  No step
extracts a divisor, so the resulting morphism \(q_Z:Z_0\to Z\) is small.
It is projective, hence \(Z_0\) is compact K\"ahler.  Codimension-one
comparison and reflexive extension identify its perturbed adjoint with
the actual pullback from \(Z\).

The same smallness gives the crepant actual pullback for the full boundary.
Moreover \(q_Z\) is an isomorphism over a smooth germ of \(Z\).  To see
this, take a relatively very ample line over a small Stein neighborhood
of the germ.  It has a meromorphic section there: after large relative
ample twists, relative Serre generation and Cartan generation supply
sections whose quotient is such a section.  Its divisor pushes to a
Cartier divisor on the smooth factorial target.  Smallness leaves no
exceptional prime, so the original divisor and line are its pullback.
A pulled-back line cannot be relatively ample on a positive-dimensional
fiber.  Any nonisomorphism fiber here would have positive dimension:
otherwise properness and the fiber-dimension theorem make the morphism
finite near that fiber, and a finite bimeromorphic morphism to the normal
target is an isomorphism.  Thus the small model is unchanged over this
smooth germ.  The lc centers of the full pair meet the
smooth crossing locus established in
Proposition~\ref{prop:strata-adjunction}; the small model is generically
unchanged there.  The full transformed pair \((Z_0,B_0)\) is therefore
dlt, and its floor \(C_0\) is the strict transform of \(C_Z\).

For a three-dimensional stratum, apply Das--Ou's lc threefold abundance
theorem \cite[Corollary~1.3]{DasOu2025} to this full pair.  Its adjoint is
the analytically nef actual pullback of the ambient restriction: the
metric lower bound restricts to \(Z\) and then pulls back.  The detailed
conventions in \cite[Section~2.1]{DasOu2025} use the reflexive canonical
sheaf, the invertible reflexive pluriadjoint, actual restriction
isomorphisms, and smooth-potential nefness.  Thus this application does not
choose a global canonical Weil divisor.  Generation descends to \(Z\) by
normality and projection formula.  Lower strata inherit generation by
restriction from a three-dimensional stratum containing them.

Choose a common sufficiently divisible even integer \(q\) for the finite
list of strata and set
\[
 L=\OO_V(q(K_V+B)),\qquad
 L_Z=L|_Z=\OO_Z(q(K_Z+B_Z)).
\]
Enlarge \(q\) so all \(L_Z\) in dimension at most three are generated.
Their complete systems followed by Stein factorization give
\begin{equation}\label{eq:stratum-semiample-map}
 f_Z:Z\longrightarrow Y_Z,\qquad L_Z\simeq f_Z^*N_Z,
\end{equation}
where \(f_Z\) has connected fibers, \(Y_Z\) is normal projective, and
\(N_Z\) is ample.  Indeed the image in projective space is projective by
Chow, and the finite analytic Stein space is projective by algebraizing
its finite coherent algebra with GAGA.  Its line is the pullback of the
tautological ample line.  In particular, for every \(k\geq0\), all
sections of \(L_Z^k\) come from \(N_Z^k\).

\subsection{A torsion-free obstruction to restriction}

The following lemma is what lets general-fiber matching control all
parameters, including those where the floor has a singular or nonreduced
scheme fiber.  Its hypothesis that the pair is klt off the floor is
essential to the canonical character calculation.

\begin{lemma}\label{lem:floor-torsion-free}
Let \((T,G)\) be a normal irreducible compact K\"ahler lc pair with effective rational
boundary, and put \(C=\floor G\) with its reduced structure.  Assume that
the pair is klt on \(T\setminus C\).  Let \(f:T\to P\) be a surjective
holomorphic map with connected fibers to a normal irreducible compact complex space.
Suppose an actual positive adjoint multiple is pulled back from a line on
\(P\).  Then
\[
 R^1f_*\OO_T(-C)\quad\text{is torsion-free on }P.
\]
Here \(\OO_T(-C)\) is the ideal of the reduced floor.
\end{lemma}
\begin{proof}
The equality with the ideal holds because on the normal \(T\) both sheaves
consist of holomorphic functions vanishing at every height-one prime in
\(C\).  Work on a small connected base open trivializing the line whose
pullback is \(m(K_T+G)\).  Its nonvanishing pulled-back frame is a local
meromorphic \(m\)-pluricanonical form \(\theta\).

Adjoin the full set of \(m\)th roots of this form.  This construction is
local even without a canonical Cartier frame.  At a normal local domain
\(A\), choose a meromorphic canonical generator \(\eta\) and write
\(\theta/\eta^m=a/b\).  The monic algebra
\[
 A[w]/(w^m-a b^{m-1})\ \subset\
 \operatorname{Frac}(A)[t]/(t^m-a/b),\qquad w=bt,
\]
is finite and reduced; its analytic normalization is finite
\cite[Part~B, Section~4, Corollaries~2--3]{Houzel1961}.  Changes of meromorphic canonical generator multiply
the root coefficient by an \(m\)th power, and identify the total algebras
and their integral closures.  They therefore glue to the normalized cover.
Keep every component and the full \(\mu_m\)-action.  The tautological
form \(\tau=t\eta\) is a meromorphic top form on a projective resolution
with smooth source; take this resolution functorially and equivariantly, after
shrinking near the compact fibers.  Write \(l:H\to T\) for the composite.

At a general prime of \(T\) with boundary coefficient \(b\), the order of
\(\tau\) upstairs is
\begin{equation}\label{eq:canonical-root-order}
 e(1-b)-1,
\end{equation}
where \(e\) is the ramification index.  This follows from the pole order
\(-eb\) of the root coefficient and the differential ramification order
\(e-1\).  It is \(-1\) when \(b=1\) and is an integer between zero and
\(e-1\) when \(0\leq b<1\).  A meromorphic form in the \(\tau\)-character
has the form \(h\tau\), where \(h\) is an invariant meromorphic function
and hence descends to \(T\); this uses the entire root algebra, whose
invariants are the normal base.  Regularity in codimension one forces
\(h\) to be holomorphic and to vanish on \(C\).

Conversely, suppose \(h\in\OO_T(-C)\).  On a log resolution of the pair,
every crepant coefficient is at most one.  A coefficient-one exceptional
prime has center in \(C\), because the pair is klt away from \(C\); the
pullback of \(h\) has positive order there.  Thus in SNC coordinates
the density \(|h|^2|\theta|^{2/m}\) has every exponent strictly above
the integrability threshold.  It is locally integrable.  Change of
variables makes \(l^*h\tau\) locally square integrable on the smooth
resolution, and a square-integrable meromorphic top form is holomorphic.
This proves the exact character equality
\begin{equation}\label{eq:canonical-character}
 (l_*\omega_H)_{\tau}=\OO_T(-C)\tau.
\end{equation}

Each component of the full normalized cover dominates the normal base,
and each resolved component does so as well.  The inverse image in \(T\)
of the connected base open is connected, because \(f\) is proper with
connected fibers; normality makes it irreducible.  Finite maps followed by
projective resolutions have K\"ahler sources near the whole inverse image
of a sufficiently small base neighborhood: patch relative ample metrics
and add a large pulled-back K\"ahler form.  Canonical torsion-freeness
\cite[Theorem~2.9 and Proposition~2.11]{Fujino2023AnalyticVanishing}
applies componentwise.  For the generically finite map to \(T\), its
higher canonical images vanish generically and hence vanish.  For the
map to \(P\), its first canonical image is torsion-free.  Leray and the
character projector in \eqref{eq:canonical-character} make
\(R^1f_*\OO_T(-C)\) a direct summand of that torsion-free sheaf.  This
proves the assertion locally, and hence on \(P\).
\end{proof}

Here is its restriction consequence.  Suppose \(L_T=f^*N\) in the lemma,
where \(P\) is projective and \(N\) ample.  The ideal sequence gives
\begin{equation}\label{eq:restriction-cokernel}
 \mathscr Q:=\coker\bigl(\OO_P\longrightarrow f_*\OO_C\bigr)
 \lhook\joinrel\longrightarrow R^1f_*\OO_T(-C).
\end{equation}
If \(C\) is vertical, meaning that it does not dominate \(P\), then
\(\mathscr Q\) is supported on a proper analytic subset and is zero by
torsion-freeness.  Serre vanishing for the kernel of
\(\OO_P\to f_*\OO_C\), after tensoring by \(N^k\), shows that every
section of \(L_T^k|_C\) extends to \(T\) for all sufficiently large \(k\).
The bound is uniform over the sections.

If \(C\) dominates \(P\), the left arrow into \(f_*\OO_C\) is injective,
and \(f_*\OO_C\) is torsion-free by \eqref{eq:restriction-cokernel}.
The finite Stein space of the reduced \(C\to P\) is reduced and every
irreducible component dominates \(P\): a vertical minimal prime of a
finite reduced algebra over a domain would supply a nonzero torsion
element.  Off a proper analytic subset this finite map is \'etale.  A
section of \(L_T^k|_C\) whose values agree throughout the reduced floor
fiber over every point of some dense open of \(P\) has zero image in
\(\mathscr Q\otimes N^k\) on that open,
and hence everywhere by torsion-freeness.  Its local lifts from \(N^k\)
are unique and glue.  Thus it extends.  This argument neither asserts
that special scheme fibers are reduced nor uses cohomology base change at
a special point.

\subsection{A Mori contraction over the semiample target}

The torsion-free argument has settled extension when the floor is
vertical over the semiample target.  For a horizontal floor, it remains
to make a section take the same value on the connected components of a
general floor fiber.  We first construct a Mori model on which those
components can be counted.

Fix a stratum \(Z\) of dimension one, two, or three and its small model
\(q_Z:(Z_0,B_0)\to(Z,B_Z)\).  Put \(C_0=\floor{B_0}\) and
\(f_0:Z_0\to Y=Y_Z\).  The full adjoint is the actual pullback of the
semiample line on \(Y\).

\begin{lemma}[A Mori model for a horizontal floor]\label{lem:horizontal-mori-model}
Assume that \(C_0\) dominates \(Y\).  There is a finite sequence of
projective divisorial contractions and flips from \(Z_0\) to an ordinary
\(\Q\)-factorial compact K\"ahler model \(T\), followed by a projective
Mori contraction
\[
 u:T\longrightarrow W,\qquad g:W\longrightarrow Y,
\]
where \(W\) is normal compact K\"ahler and both maps have connected
fibers.  Write \(G\) for the transform of \(B_0\) and
\(C^+=\floor G\).  The pair \((T,G)\) is effective lc and klt off
\(C^+\); the divisor \(C^+\) is relatively ample for \(u\) and still
dominates \(Y\).  All steps factor over \(Y\) and extract no divisor.
On a common projective resolution of the sequence, the full adjoints
are equal as meromorphic pullbacks of the same actual line from \(Y\).
\end{lemma}
\begin{proof}

Choose a small rational \(\varepsilon>0\) for which
\(B_0-\varepsilon C_0\) is effective klt.  This uses dlt and the fact
that all its lc centers are contained in the floor.  Let \(H\) be an
integral very ample divisor on \(Y\) and choose a rational effective
\(H'\sim_\Q dH\) with \(d>2\dim Z\) such that
\[
 B_0-\varepsilon C_0+f_0^*H'
\]
is still klt.  A sum of sufficiently many general members of \(|H|\)
with small rational weights does this by Bertini on a log resolution.
If \(Y\) is a point take \(H'=0\).

The adjoint of this perturbed pair is not pseudo-effective.  On a smooth
general fiber of a resolution of \(f_0\), its class is the negative of
the nonzero effective divisor \(\varepsilon C_0\) restricted to that
fiber: the full adjoint and \(H'\) are pulled back from \(Y\), and the
floor is horizontal.  Its integral against a K\"ahler power is strictly
negative.  If the ambient class were pseudo-effective, the potential of
a positive current representing its pullback would restrict to a positive
current on almost every such smooth fiber, contradicting this integral.
The smooth horizontal-fiber conditions hold on a nonempty open, so the
almost-everywhere restriction is sufficient.

This same integral also produces a negative class in the cone used by the
K\"ahler cone theorem.  If the smooth fiber has dimension \(v>0\), push
its positive closed bidimension-\((1,1)\) current \(\omega_F^{v-1}\)
through the fiber inclusion and the resolution to \(Z_0\).  It represents
a class in \(\overline{\operatorname{NA}}(Z_0)\), and its pairing with
the perturbed adjoint is exactly that strictly negative integral.  The
cone decomposition therefore supplies a negative extremal ray.

Run the klt K\"ahler program for this perturbed adjoint.  In dimension
three, the negative extremal-ray and contraction theorems
\cite[Corollary~5.3 and Theorem~5.5]{DasHaconPaun2024} are used through
Subsection~\ref{subsec:threefold-contraction-bridge}, which incorporates
Proposition~\ref{prop:prime-floor-extension} and the specified external
results.  They give
connected projective contractions to normal compact K\"ahler targets.
The supporting class differs from the klt adjoint by a K\"ahler class,
which is the projectivity condition in that theorem.  The supporting
class need not be big, and the perturbed adjoint here is not
pseudo-effective.  Flips and termination
of flip sequences are provided by
\cite[Theorems~2.24--2.25]{DasHacon2024}.  In dimensions at most two the
non-pseudo-effective pair is projective and the ordinary projective program
applies.  For a surface, otherwise a non-Moishezon K\"ahler resolution
has a nonzero holomorphic two-form: if \(H^{2,0}=0\), rational approximation
of a K\"ahler class would make it projective.  That form descends
reflexively and, with the effective boundary, gives a section of a positive
adjoint multiple, contradicting non-pseudo-effectivity.  A Moishezon
K\"ahler klt surface is projective by Namikawa's criterion
\cite[Corollary~6']{Namikawa2002}.  Curves are projective.

We specify the ordinary \(\Q\)-factorial and finiteness details used by
this program.  For a projective negative extremal contraction, a global
rational line of degree zero on all contracted curves descends in a
multiple by \cite[Theorem~2.44(3c)]{DasHaconPaun2024}.  Here we apply its
projective relative theorem to the effective klt pair, with the compact
set equal to the entire compact target.  For this projective surjection,
property~Q of \cite[Definition~2.42]{DasHaconPaun2024} requires a normal
source and compact source and target, as here.  All fiber curves span the one relative
ray, so the theorem's contraction is the given contraction.  Indeed its
map is constant on each connected projective fiber of the given contraction,
and rigidity factors it through that contraction.  Its structure map back
to the relative base supplies the reverse factorization; surjectivity makes
the two factorizations inverse.  This use of Theorem~2.44 does
not assert a factoriality hypothesis for that theorem.  At a divisorial step, adjust the
transform of a destination Weil divisor by a \(\Q\)-Cartier exceptional
prime of nonzero degree on the ray, and descend the resulting degree-zero
line.  Comparison in codimension one makes the destination divisor
\(\Q\)-Cartier.  At a flip, adjust on the negative side by the negative
adjoint, descend, and pull to the positive side; the positive adjoint is
\(\Q\)-Cartier and the same comparison applies.  Boundary components
and the adjoint then also give the canonical reflexive power.  This
preserves ordinary \(\Q\)-factoriality.

The rank of global rational divisor classes modulo curve numerical
equivalence is finite.  Each running model is a connected compact complex
analytic space.  Its reduced underlying real analytic space has bounded
Zariski tangent dimension by a finite analytic chart cover.
Acquistapace--Broglia--Tognoli \cite[Theorem~1]{AcquistapaceBrogliaTognoli1979}
embed it closedly in Euclidean space, and \L{}ojasiewicz
\cite[Theorem~1]{Lojasiewicz1964} gives a compatible locally finite
triangulation; the subcomplex corresponding to the compact model is
finite.  Hence its rational \(H^2\) is finite-dimensional.  The first
Chern class sends global rational Cartier divisors to that group.  Every
divisor in its kernel has degree zero on each compact curve, by restriction
to the curve's normalization.  The rational divisor space modulo curve
numerical equivalence is consequently a quotient of the Chern-class image,
and has finite rank.  The rank drops at a divisorial contraction, by the
nonzero exceptional ray degree, and does not increase at a flip.  For
the latter assertion a numerically trivial line on the negative side
descends as above; the descended line is curve-numerically trivial because
every curve downstairs is covered by a curve upstairs, using projectivity
over that curve.  Its pullback is trivial on curves on the positive side,
and every divisor there is a transform by smallness.  There can therefore
be only finitely many divisorial steps.  Together with termination of flip
sequences this gives a finite program to a nef model or a Mori contraction.
All steps extract no divisors.

Every step factors over \(Y\).  Inductively write the running perturbed
adjoint as
\[
 A_i^{\varepsilon}=K_{X_i}+B_i-\varepsilon C_i+f_i^*H'.
\]
For a negative contraction apply the relative cone and length theorem
\cite[Theorem~2.45]{DasHaconPaun2024} to \(K_{X_i}+B_i-\varepsilon C_i\)
over its entire compact contraction target.  The source is ordinary
\(\Q\)-factorial and the pair is klt.  Property~Q holds as above; because
the compact set is the whole target, Theorem~2.45's factoriality over
that set is precisely the global factoriality of the source, including
its canonical reflexive power.  A negative generator \(\Gamma\) has
\[
 -(K_{X_i}+B_i-\varepsilon C_i)\cdot\Gamma\leq2\dim Z.
\]
If \(f_i(\Gamma)\) were a curve, then \(f_i^*H'\cdot\Gamma\geq d\),
contradicting the negativity of \(A_i^{\varepsilon}\).  On the nonnegative
part of that relative cone, both \(K_{X_i}+B_i-\varepsilon C_i\) and the nef
\(f_i^*H'\) are nonnegative, while their sum is nonpositive on the
contracted cone; its \(H'\)-degree is zero there as well.  All contracted
curves are vertical over \(Y\).  A holomorphic map from a connected
projective fiber that is nonconstant would map a curve nontrivially, so
\(f_i\) is constant on each fiber.  Rigidity and normality of the target
give the factorization, and the positive side of a flip factors through
the same base.

We check the boundary properties needed to repeat the step.  No prime is
extracted, so the transformed full boundary is effective and its floor
is exactly the transform \(C_{i+1}\) of \(C_i\).  The running MMP
preserves the klt pair
\(B_i-\varepsilon C_i+f_i^*H'\).  Removing the effective divisor
\(f_{i+1}^*H'\) shows that \(B_{i+1}-\varepsilon C_{i+1}\) is klt.
In particular the full pair is klt away from its floor.

The full adjoints remain actually crepant pullbacks from \(Y\).
Indeed the chosen meromorphic pluriadjoint identity extends to the new
model in codimension one, since no prime is extracted, and then
reflexively.  On a common graph both meromorphic maps start from the same
line pulled back from \(Y\) and agree on the dense isomorphism open.
They agree everywhere as meromorphic maps.  Equality of discrepancies
therefore preserves log canonicity of the full pair.  If \(r,s\) are
the graph projections, negativity for the perturbed step gives
\begin{equation}\label{eq:horizontal-persistence}
 0\leq r^*A_i^{\varepsilon}-s^*A_{i+1}^{\varepsilon}
   =\varepsilon(s^*C_{i+1}-r^*C_i).
\end{equation}
A horizontal component of \(C_i\) cannot disappear into an entirely
vertical \(C_{i+1}\), since the coefficient of the right side along its
strict valuation would then be negative.  Thus the floor remains horizontal.
On every subsequent model the full adjoint and \(H'\) still come from
\(Y\).  Repeating the resolved smooth-fiber integral and pushing its
K\"ahler-power current as above gives a negative class in its
\(\overline{\operatorname{NA}}\) cone and hence a negative extremal ray.
In particular a nef endpoint is impossible.

We have obtained a Mori contraction
\[
 u:T\longrightarrow W,\qquad g:W\longrightarrow Y,
\]
with \(T\) normal compact K\"ahler, \(u\) projective of fiber type, and
both maps having connected fibers.  For \(g\), this follows also from
\((g u)_*\OO_T=\OO_Y\) and \(u_*\OO_T=\OO_W\).  The full boundary
\(G\) is effective lc, its adjoint is an actual pullback from \(Y\), and
it is klt away from \(C^+=\floor G\): subtracting \(\varepsilon C^+\)
is klt.  The running perturbed adjoint satisfies
\[
 A^{\varepsilon}\equiv_u-\varepsilon C^+,
 \qquad -A^{\varepsilon}\equiv_u+\varepsilon C^+.
\]
The left side of the second equivalence is relatively ample for the
Mori contraction.  Thus \(C^+\) is relatively ample for \(u\).
The graphs of the divisorial steps and flips are projective over both
sides.  A main component of their iterated fiber product, followed by
normalization and projective resolution, gives the common projective
resolution in the statement.  Its smooth source is compact K\"ahler.
\end{proof}

\subsection{The restriction criterion}

We now use the Mori model only to compare the values of a boundary
section on general fibers.  The torsion-free obstruction then extends
that comparison over every parameter.

\begin{proposition}[Restriction criterion]\label{prop:restriction-link}
For every stratum \(Z\) of dimension one, two, or three, there is either
no comparison or one proper bimeromorphic comparison between two normal
components of \(C_Z\), allowing a component to be compared with itself,
with the following properties.
\begin{enumerate}
 \item On a resolution of its graph with smooth source, the two pulled-back actual
 boundary adjoint lines are equal as invertible subsheaves of meromorphic
 pluricanonical forms.
 \item For all sufficiently large divisible \(k\), uniformly over
 sections, a section of \(L_Z^k|_{C_Z}\) extends to \(Z\) if its two
 pullbacks agree under this comparison.  With no comparison there is no
 matching condition.
 \item When present, the comparison pairs two coefficient-one markings
 on general projective-line fibers of a projective Mori contraction on
 a bimeromorphic compact K\"ahler model of \(Z\).
\end{enumerate}
\end{proposition}
\begin{proof}
If \(C_Z\) is vertical over \(Y_Z\), the consequence of
\eqref{eq:restriction-cokernel} gives extension in a uniform large-degree
tail without a comparison.  Suppose therefore that it is horizontal,
and use the small model \(q_Z\) and the Mori model
\(u:T\to W\), \(g:W\to Y\) of Lemma~\ref{lem:horizontal-mori-model}.
Extending the pulled-back section on \(C_0\) suffices: a section on
\(Z_0\) descends to \(Z\), and equality of its restriction after
pullback implies equality on the reduced \(C_Z\), since every component
is covered birationally.

\subsection*{The floor on general Mori fibers}

We now compare the connected components of a general floor fiber with
the markings of one Mori fiber.  On a common projective log resolution
of the program, the full crepant boundary is the same on both sides.
Its coefficient-one union maps to both \(C_0\) and \(C^+\) with connected
fibers by Lemma~\ref{lem:floor-connectedness}.  After restriction over
any \(y\in Y\), proper closed maps with connected fibers preserve
connected components.  Thus the connected components of the underlying
spaces of \(C_{0,y}\) and \(C^+_y\) correspond.  This is a topological
statement, not a reducedness claim for their scheme fibers.

The finite Stein space of \(C^+\to W\) has every component dominating
\(W\).  Indeed relative ampleness makes \(C^+\to W\) surjective, and
Lemma~\ref{lem:floor-torsion-free} and
\eqref{eq:restriction-cokernel} give the torsion-free property used in
the preceding discussion.  For general \(y\), every irreducible component
of the finite Stein fiber has dimension \(\dim W-\dim Y\), by the
dimension theorem.  The fiber \(W_y\) is irreducible of that dimension:
a resolution of \(W\) has connected fibers over \(Y\), and a general
one is smooth and connected, hence irreducible, and surjects onto \(W_y\).
Each component of the finite Stein fiber therefore dominates \(W_y\).
It follows that every connected component of \(C^+_y\) meets
\(u^{-1}(w)\) for a common general \(w\in W_y\).

The underlying space of a general fiber \(F\) of \(u\) is irreducible
by the same resolution argument.  If \(\dim F\geq2\), a Cartier multiple
of \(C^+\) restricts to a nonzero effective ample Cartier divisor on
\(F_{\red}\).  Its support is connected.  To recall the reason, if an
ample effective divisor split into two disjoint nonzero parts, general
hyperplane sections would reduce to an irreducible projective surface
\(F_2\).  Write the two induced nonzero effective Cartier parts on it as
\(D_1,D_2\), and let \(v:\widetilde F\to F_2\) be a projective resolution
with smooth source.  Their pullbacks are effective and orthogonal.  Although
\(v^*(D_1+D_2)\) is only nef and big, the projection formula gives
\[
 v^*(D_1+D_2)\cdot v^*D_i
 =(D_1+D_2)\cdot v_*[v^*D_i]>0\qquad(i=1,2):
\]
the pushforward is a nonzero effective curve cycle and \(D_1+D_2\) is
ample on the integral surface.  Orthogonality then gives
\((v^*D_i)^2>0\) for both \(i\).  The surface Hodge index theorem forbids
two such orthogonal positive classes.  This argument does not require
normality of \(F_{\red}\).

If \(\dim F=1\), choose \(F\) also in the generic-smoothness open.
Generality avoids the singular locus of \(T\), the
singularities and intersections of the boundary, and the ramification of
its horizontal primes.  The fiber is a smooth connected curve, and the
full pullback identity gives
\[
 0=\deg(K_F+G|_F)=2g(F)-2+\deg(G|_F).
\]
There is at least one coefficient-one point.  Hence \(F\simeq\PP^1\),
and its floor has one or two points.  In the two-point case these points
exhaust the horizontal boundary on \(F\).  In all connected cases,
every connected component of \(C^+_y\) meets the same connected subset
\(C^+\cap F\), so \(C^+_y\), and therefore \(C_{0,y}\), is connected.
Then no comparison is needed.  In the two-point case each connected
component of \(C_{0,y}\) is represented by at least one of these two
markings, after their transfer through the common graph.

\subsection*{The two-marking comparison}

Normalize each surviving horizontal prime \(D\subset C^+\) and take the
Stein factorization
\[
 D^\nu\longrightarrow E_D\longrightarrow W.
\]
The first map is projective bimeromorphic and \(E_D\) is normal; the
second is finite.  If there are two horizontal primes, each finite map
has degree one and is an isomorphism over the normal \(W\).  The main
component of the fiber product of their normalizations over \(W\) gives
a proper bimeromorphic comparison.  If there is one horizontal prime,
\(E_D\to W\) has degree two.  On a dense open it is \'etale and has an
exchange.  The reduced horizontal non-diagonal component of
\(E_D\times_W E_D\) has finite generically one-to-one projections to
the normal \(E_D\); both are isomorphisms.  It defines the exchange
involution globally, including across the branch.  The main component
of \(D^\nu\times_{E_D,\mathrm{exchange}}D^\nu\) lifts it to a proper
bimeromorphic graph.  It is the Stein space of the normalized prime
that is used here, not the possibly nonnormal Stein space of the entire
floor.

Compose this graph with the common graphs of the program and with the
small model.  It gives the asserted comparison between normal primes of
\(C_Z\), possibly a self-comparison.  Each surviving prime is generically
finite over \(W\); the image in \(W\) of the proper subset where its
transfer is not an isomorphism is proper.  Thus both markings on a general
Mori fiber lie in the common transfer isomorphism loci.

These graph projections are projective.  The maps \(D^\nu\to E_D\)
are projective and the maps \(E_D\to W\) are finite; their fiber products
and closed main components are therefore projective over the branches.
The small-model and MMP graphs have the same property.  Iterated main
components and finite normalization preserve it, so the comparison admits
a projective resolution with smooth compact K\"ahler source.

We verify the actual meromorphic identity required by property (1).
The two residues of a logarithmic fiber form give the same comparison
as the even Poincar\'e-residue diagrams of
\cite[Section~3, Definition~13 and Proposition~14]{Kollar2013Sources};
we include the calculation for the normalized two-branch construction
above.
On a smooth ruled open of \(W\), order the two markings after an \'etale
local cover when necessary, choose a base volume form \(\xi\), and a
fiber coordinate with markings at zero and infinity.  The full
\(q\)-pluriadjoint frame coming from \(Y\) has the form
\[
 a(w)(dz/z)^{\otimes q}\otimes\xi^{\otimes q}.
\]
There is no other horizontal divisor on the general fiber.  Its two
residues are \(a(w)\xi^{\otimes q}\) and
\((-1)^q a(w)\xi^{\otimes q}\); they agree because \(q\) is even.
The calculation is invariant under exchanging the two local markings.
Full crepancy transfers it to the original primes.  On a resolution of
their comparison graph, both actual adjoint lines are the same pullback
of \(N_Z\) from \(Y\).  Their meromorphic embeddings agree on this dense
ruled open, hence agree everywhere.  Thus the invertible subsheaves of
meromorphic pluriforms are equal even at exceptional primes.  An abstract
Q-linear equivalence without this residue calculation would not suffice.

A section satisfying this comparison has equal values at the two
representatives over a general \(w\).  On every connected component of
a compact reduced floor fiber it is constant in the trivialized pulled-back
line: holomorphic functions on a compact connected reduced complex space
are constant.  The representatives meet every component, so its values
agree on \(C_{0,y}\) for every \(y\) in a dense base open.  The horizontal consequence of
\eqref{eq:restriction-cokernel} extends it to \(Z_0\), and then it descends
to \(Z\).  The only asymptotic bound on \(k\) came from Serre vanishing in
the vertical case.  Finitely many strata admit a common sufficiently
divisible tail.  This proves all three properties.
\end{proof}

For the remainder of the boundary proof, call each comparison supplied
by Proposition~\ref{prop:restriction-link} a \emph{link}.  This word
refers to its proper bimeromorphic graph together with the proved equality
of meromorphic adjoint subsheaves.  It does not mean only an abstract
linear equivalence.  We will use all curve comparisons, but on surfaces
only the groupoid generated by links of three-dimensional strata.  That
restriction is what allows the common ambient K\"ahler class to control
their scalar action.
